% Distilled blueprint for Evans, Appendix D (Linear Functional Analysis).

\section{Linear Functional Analysis}
\label{sec:linear-functional-analysis}

\subsection{Banach Spaces}
\label{sec:banach-spaces}

Let $X$ be a real linear space.

\begin{definition}[Norm]
\label{def:norm-normed-linear-space}
\dcref{appD:D.1-def-1}
\group{Functional Analysis}
A map $\|\cdot\|:X\to[0,\infty)$ is a \textit{norm} if, for $u,v\in X$ and $\lambda\in\mathbb R$,

(i) $\|u+v\|\leq\|u\|+\|v\|$;

(ii) $\|\lambda u\|=|\lambda|\|u\|$;

(iii) $\|u\|=0$ if and only if $u=0$.

Condition (i) is the triangle inequality.
\end{definition}

Hereafter $X$ is a normed real linear space.

\begin{definition}[Convergence in a normed space]
\label{def:convergence-normed-space}
\uses{def:norm-normed-linear-space}
\dcref{appD:D.1-def-2}
\group{Functional Analysis}

A sequence $\{u_k\}_{k=1}^\infty\subset X$ converges to $u\in X$, written $u_k\to u$, if
\[
\lim_{k\to\infty}\|u_k-u\|=0.
\]
\end{definition}

\begin{definition}[Cauchy sequence, completeness, and Banach space]
\label{def:cauchy-sequence-complete-banach-space}
\uses{def:norm-normed-linear-space}
\dcref{appD:D.1-def-3}
\group{Functional Analysis}

(i) A sequence $\{u_k\}_{k=1}^\infty\subset X$ is \textit{Cauchy} if, for every $\varepsilon>0$, there is $N$ such that
\[
\|u_k-u_l\|<\varepsilon\qquad(k,l\geq N).
\]

(ii) The space $X$ is \textit{complete} if every Cauchy sequence converges to an element of $X$.

(iii) A \textit{Banach space} is a complete normed linear space.
\end{definition}

\begin{definition}[Separable space]
\label{def:separable-space}
\dcref{appD:D.1-def-4}
\group{Functional Analysis}
The space $X$ is \textit{separable} if it contains a countable dense subset.
\end{definition}

\begin{example}[Banach space examples: $L^p$, H\"older, and Sobolev spaces]
\label{ex:banach-space-examples-lp-holder-sobolev}
\uses{def:cauchy-sequence-complete-banach-space, def:holder-space, def:sobolev-space-wkp}
\dcref{appD:D.1-ex-1}
\group{Functional Analysis}

(i) If $U\subset\mathbb R^n$ is open and $1\leq p\leq\infty$, set
\[
\|f\|_{L^p(U)}:=
\begin{cases}
\left(\int_U|f|^p\,dx\right)^{1/p} & 1\leq p<\infty,\\
\operatorname*{ess\,sup}_U|f| & p=\infty.
\end{cases}
\]
The measurable functions with finite norm form the Banach space $L^p(U)$ after functions agreeing almost everywhere are identified.

(ii) H\"older spaces are defined in \cref{def:holder-space}.

(iii) Sobolev spaces are defined in \cref{def:sobolev-space-wkp}.
\end{example}

\subsection{Hilbert Spaces}
\label{sec:hilbert-spaces}

Let $H$ be a real linear space.

\begin{definition}[Inner product]
\label{def:inner-product}
\dcref{appD:D.2-def-1}
\group{Functional Analysis}
A map $(\cdot,\cdot):H\times H\to\mathbb R$ is an \textit{inner product} if:

(i) $(u,v)=(v,u)$ for all $u,v\in H$;

(ii) $u\mapsto(u,v)$ is linear for each $v\in H$;

(iii) $(u,u)\geq0$ for every $u\in H$;

(iv) $(u,u)=0$ if and only if $u=0$.
\end{definition}

\begin{definition}[Norm induced by an inner product]
\label{def:norm-from-inner-product}
\uses{def:inner-product}
\dcref{appD:D.2-def-2}
\group{Functional Analysis}

The norm induced by an inner product is
\[
(1)\qquad \|u\|:=(u,u)^{1/2}\qquad(u\in H).
\]
\end{definition}

\begin{lemma}[Cauchy--Schwarz inequality in a Hilbert space]
\label{lem:cauchy-schwarz-inequality-hilbert}
\uses{lem:cauchy-schwarz-inequality, def:inner-product, def:norm-from-inner-product}
\dcref{appD:D.2-lem-1}
\group{Functional Analysis}

For $u,v\in H$,
\[
(2)\qquad |(u,v)|\leq\|u\|\|v\|.
\]
This is the inner-product form of \cref{lem:cauchy-schwarz-inequality}.
\end{lemma}

The Cauchy--Schwarz inequality and the inner-product axioms imply that (1) is a norm.

\begin{definition}[Hilbert space]
\label{def:hilbert-space}
\uses{def:cauchy-sequence-complete-banach-space, def:inner-product, def:norm-from-inner-product}
\dcref{appD:D.2-def-3}
\group{Functional Analysis}

A \textit{Hilbert space} is a Banach space whose norm is induced by an inner product.
\end{definition}

\begin{example}[$L^2(U)$ is a Hilbert space]
\label{ex:l2-hilbert-space-example}
\uses{def:hilbert-space, def:inner-product}
\dcref{appD:D.2-ex-1}
\group{Functional Analysis}

The space $L^2(U)$ is a Hilbert space with
\[
(f,g)=\int_Ufg\,dx.
\]
\end{example}

\begin{example}[$H^1(U)$ is a Hilbert space]
\label{ex:h1-hilbert-space-example}
\uses{def:hilbert-space, def:inner-product}
\dcref{appD:D.2-ex-2}
\group{Functional Analysis}

The Sobolev space $H^1(U)$ is a Hilbert space with
\[
(f,g)=\int_U(fg+Df\cdot Dg)\,dx.
\]
\end{example}

\begin{definition}[Orthogonality and orthonormal bases]
\label{def:orthogonal-orthonormal-basis}
\uses{def:hilbert-space}
\dcref{appD:D.2-def-4}
\group{Functional Analysis}

(i) Elements $u,v\in H$ are \textit{orthogonal} if $(u,v)=0$.

(ii) A countable basis $\{w_k\}_{k=1}^\infty\subset H$ is \textit{orthonormal} if
\[
(w_k,w_l)=0\quad(k\ne l),
\qquad
\|w_k\|=1\quad(k=1,2,\ldots).
\]
For every $u\in H$,
\[
u=\sum_{k=1}^\infty(u,w_k)w_k
\quad\text{in }H,
\qquad
\|u\|^2=\sum_{k=1}^\infty(u,w_k)^2.
\]
\end{definition}

\begin{definition}[Orthogonal complement]
\label{def:orthogonal-complement-subspace}
\uses{def:orthogonal-orthonormal-basis}
\dcref{appD:D.2-def-5}
\group{Functional Analysis}

For a subspace $S\subset H$,
\[
S^\perp:=\{u\in H\mid(u,v)=0\text{ for every }v\in S\}.
\]
\end{definition}

\subsection{Bounded Linear Operators}
\label{sec:bounded-linear-operators}

Let $X$ and $Y$ be real Banach spaces.

\begin{definition}[Linear operator]
\label{def:linear-operator-banach}
\uses{def:cauchy-sequence-complete-banach-space}
\dcref{appD:D.3-def-1}
\group{Functional Analysis}

(i) A map $A:X\to Y$ is a \textit{linear operator} if
\[
A[\lambda u+\mu v]=\lambda Au+\mu Av
\]
for all $u,v\in X$ and $\lambda,\mu\in\mathbb R$.

(ii) Its range and null space are
\[
R(A):=\{v\in Y\mid v=Au\text{ for some }u\in X\},
\qquad
N(A):=\{u\in X\mid Au=0\}.
\]
\end{definition}

\begin{definition}[Bounded linear operator]
\label{def:bounded-linear-operator}
\uses{def:linear-operator-banach}
\dcref{appD:D.3-def-2}
\group{Functional Analysis}

A linear operator $A:X\to Y$ is \textit{bounded} if
\[
\|A\|:=\sup\{\|Au\|_Y\mid\|u\|_X\leq1\}<\infty.
\]
\end{definition}

Every bounded linear operator $A:X\to Y$ is continuous.

\begin{definition}[Closed linear operator]
\label{def:closed-linear-operator}
\dcref{appD:D.3-def-3}
\group{Functional Analysis}
A linear operator $A:X\to Y$ is \textit{closed} if
\[
u_k\to u\text{ in }X,\quad Au_k\to v\text{ in }Y
\quad\Longrightarrow\quad Au=v.
\]
\end{definition}

\begin{theorem}[Closed Graph Theorem]
\label{thm:closed-graph-theorem}
\uses{def:closed-linear-operator}
\dcref{appD:thm-1}
\group{Functional Analysis}

Every closed linear operator $A:X\to Y$ is bounded.
\end{theorem}

\begin{definition}[Resolvent set and spectrum]
\label{def:resolvent-set-spectrum-appendix-d}
\dcref{appD:D.3-def-4}
\group{Functional Analysis}
Let $A:X\to X$ be bounded and linear.

(i) Its resolvent set is
\[
\rho(A):=\{\eta\in\mathbb R\mid A-\eta I
\text{ is one-to-one and onto}\}.
\]

(ii) Its spectrum is
\[
\sigma(A):=\mathbb R-\rho(A).
\]
\end{definition}

\begin{remark}[The resolvent yields a bounded inverse]
\label{rem:resolvent-inverse-bounded}
\uses{thm:closed-graph-theorem, def:resolvent-set-spectrum-appendix-d}
\dcref{appD:D.3-rem-1}
\group{Functional Analysis}

If $\eta\in\rho(A)$, then $(A-\eta I)^{-1}:X\to X$ is a bounded linear operator by the Closed Graph Theorem.
\end{remark}

\begin{definition}[Eigenvalue and eigenvector]
\label{def:eigenvalue-eigenvector-appendix-d}
\uses{def:resolvent-set-spectrum-appendix-d}
\dcref{appD:D.3-def-5}
\group{Functional Analysis}

(i) A point $\eta\in\sigma(A)$ is an \textit{eigenvalue} if
\[
N(A-\eta I)\ne\{0\}.
\]
The set of eigenvalues is the point spectrum $\sigma_p(A)$.

(ii) If $w\ne0$ and $Aw=\eta w$, then $w$ is an eigenvector associated with $\eta$.
\end{definition}

\begin{definition}[Bounded linear functional and dual space]
\label{def:bounded-linear-functional-dual-space}
\dcref{appD:D.3-def-6}
\group{Functional Analysis}
(i) A bounded linear operator $u^*:X\to\mathbb R$ is a \textit{bounded linear functional}.

(ii) The dual space $X^*$ is the space of all bounded linear functionals on $X$.
\end{definition}

\begin{definition}[Duality pairing, dual norm, and reflexivity]
\label{def:dual-pairing-reflexive-banach-space}
\uses{def:bounded-linear-functional-dual-space}
\dcref{appD:D.3-def-7}
\group{Functional Analysis}

(i) For $u\in X$ and $u^*\in X^*$, the duality pairing is
\[
\langle u^*,u\rangle:=u^*(u).
\]

(ii) The dual norm is
\[
\|u^*\|:=\sup\{\langle u^*,u\rangle\mid\|u\|\leq1\}.
\]

(iii) A Banach space is \textit{reflexive} if the canonical embedding identifies $(X^*)^*$ with $X$: for every $u^{**}\in(X^*)^*$ there is $u\in X$ such that
\[
\langle u^{**},u^*\rangle=\langle u^*,u\rangle
\qquad(u^*\in X^*).
\]
\end{definition}

Let $H$ now be a real Hilbert space with inner product $(\cdot,\cdot)$.

\begin{theorem}[Riesz Representation Theorem]
\label{thm:riesz-representation-theorem}
\uses{def:hilbert-space, def:bounded-linear-functional-dual-space}
\dcref{appD:thm-2}
\group{Functional Analysis}

For every $u^*\in H^*$ there is a unique $u\in H$ such that
\[
\langle u^*,v\rangle=(u,v)
\qquad(v\in H).
\]
The map $u^*\mapsto u$ is a linear isomorphism from $H^*$ onto $H$.
\end{theorem}

\begin{definition}[Adjoint operator and symmetric operator]
\label{def:adjoint-symmetric-operator}
\uses{thm:riesz-representation-theorem}
\dcref{appD:D.3-def-8}
\group{Functional Analysis}

(i) If $A:H\to H$ is bounded and linear, its adjoint $A^*:H\to H$ is characterized by
\[
(Au,v)=(u,A^*v)
\qquad(u,v\in H).
\]

(ii) The operator $A$ is \textit{symmetric} if $A^*=A$.
\end{definition}

\subsection{Weak Convergence}
\label{sec:weak-convergence-appendix-d}

Let $X$ be a real Banach space.

\begin{definition}[Weak convergence]
\label{def:weak-convergence}
\uses{def:bounded-linear-functional-dual-space}
\dcref{appD:D.4-def-1}
\group{Functional Analysis}

A sequence $\{u_k\}_{k=1}^\infty\subset X$ converges \textit{weakly} to $u\in X$, written $u_k\rightharpoonup u$, if
\[
\langle u^*,u_k\rangle\to\langle u^*,u\rangle
\]
for every $u^*\in X^*$.
\end{definition}

\begin{remark}[Basic properties of weak convergence]
\label{rem:weak-convergence-basic-properties}
\uses{def:weak-convergence, def:convergence-normed-space}
\dcref{appD:D.4-rem-1}
\group{Functional Analysis}

Strong convergence implies weak convergence. Every weakly convergent sequence is bounded, and
\[
u_k\rightharpoonup u
\quad\Longrightarrow\quad
\|u\|\leq\liminf_{k\to\infty}\|u_k\|.
\]
\end{remark}

\begin{theorem}[Weak compactness]
\label{thm:weak-compactness-reflexive-banach}
\uses{def:weak-convergence, def:dual-pairing-reflexive-banach-space}
\dcref{appD:thm-3}
\group{Functional Analysis}

If $X$ is reflexive and $\{u_k\}_{k=1}^\infty\subset X$ is bounded, then some subsequence satisfies
\[
u_{k_j}\rightharpoonup u
\]
for an element $u\in X$.
\end{theorem}

Consequently, every bounded sequence in a Hilbert space has a weakly convergent subsequence.

\begin{theorem}[Mazur's Theorem]
\label{thm:mazurs-theorem}
\uses{def:weak-convergence}
\dcref{appD:D.4-thm-mazur}
\group{Functional Analysis}

Every convex closed subset of $X$ is weakly closed.
\end{theorem}

\begin{example}[Weak convergence in $L^p$]
\label{ex:weak-convergence-lp-duality}
\uses{def:weak-convergence, thm:weak-compactness-reflexive-banach, def:dual-pairing-reflexive-banach-space}
\dcref{appD:D.4-ex-1}
\group{Functional Analysis}

Let $U\subset\mathbb R^n$ be open, $1\leq p<\infty$, and let $q$ satisfy
\[
\frac1p+\frac1q=1,
\qquad 1<q\leq\infty.
\]
The dual of $L^p(U)$ is $L^q(U)$: every bounded linear functional has the form
\[
f\longmapsto\int_Ugf\,dx
\qquad(g\in L^q(U)).
\]
Thus $f_k\rightharpoonup f$ weakly in $L^p(U)$ means
\[
(3)\qquad
\int_Ugf_k\,dx\longrightarrow\int_Ugf\,dx
\quad\text{for every }g\in L^q(U).
\]
The space $L^p(U)$ is reflexive for $1<p<\infty$; hence every bounded sequence in this range has a weakly convergent subsequence by \cref{thm:weak-compactness-reflexive-banach}. Convergence in (3) need not imply pointwise or almost-everywhere convergence; increasingly rapid oscillation may prevent both.
\end{example}

\subsection{Compact Operators and Fredholm Theory}
\label{sec:compact-operators-fredholm-theory}

Let $X$ and $Y$ be real Banach spaces.

\begin{definition}[Compact operator]
\label{def:compact-linear-operator}
\dcref{appD:D.5-def-1}
\group{Fredholm Theory}
A bounded linear operator $K:X\to Y$ is \textit{compact} if, for every bounded sequence $\{u_k\}\subset X$, some subsequence $\{Ku_{k_j}\}$ converges in $Y$.
\end{definition}

\begin{remark}[Compact operators map weakly convergent to strongly convergent sequences]
\label{rem:compact-operator-weak-to-strong-convergence}
\uses{def:compact-linear-operator, def:weak-convergence}
\dcref{appD:D.5-rem-1}
\group{Fredholm Theory}

If $H$ is a real Hilbert space, $K:H\to H$ is compact, and $u_k\rightharpoonup u$, then $Ku_k\to Ku$ strongly.
\end{remark}

\begin{theorem}[Compactness of adjoints]
\label{thm:compactness-of-adjoint-operator}
\uses{def:compact-linear-operator, def:adjoint-symmetric-operator}
\dcref{appD:thm-4}
\group{Fredholm Theory}

If $K:H\to H$ is compact, then $K^*:H\to H$ is compact.
\end{theorem}
\begin{proof}
\sketch Let $\{u_k\}$ be bounded and choose a subsequence $u_{k_j}\rightharpoonup u$. Since $K^*$ is bounded, $K^*(u_{k_j}-u)\rightharpoonup0$; compactness of $K$ then gives
\[
KK^*(u_{k_j}-u)\to0.
\]
Because $\{u_{k_j}-u\}$ is bounded,
\[
\|K^*(u_{k_j}-u)\|^2
=(KK^*(u_{k_j}-u),u_{k_j}-u)\to0.
\]
Thus $\{K^*u_{k_j}\}$ converges strongly.
\end{proof}

\begin{theorem}[Fredholm alternative]
\label{thm:fredholm-alternative}
\uses{thm:compactness-of-adjoint-operator, def:compact-linear-operator}
\dcref{appD:thm-5}
\group{Fredholm Theory}

Let $K:H\to H$ be compact and linear. Then:

(i) $N(I-K)$ is finite dimensional;

(ii) $R(I-K)$ is closed;

(iii) $R(I-K)=N(I-K^*)^\perp$;

(iv) $N(I-K)=\{0\}$ if and only if $R(I-K)=H$;

(v) $\dim N(I-K)=\dim N(I-K^*)$.
\end{theorem}
\begin{proof}
\sketch If $N(I-K)$ contained an infinite orthonormal set $\{u_k\}$, then $Ku_k=u_k$ and $\|Ku_k-Ku_l\|=\sqrt2$ for $k\ne l$, contradicting compactness. This proves (i).

There is $\gamma>0$ such that
\[
(4)\qquad
\|(I-K)u\|\geq\gamma\|u\|
\qquad(u\in N(I-K)^\perp).
\]
Otherwise unit vectors $u_k\in N(I-K)^\perp$ would satisfy
\[
(5)\qquad (I-K)u_k\to0.
\]
After taking a weakly convergent subsequence, compactness makes $Ku_k$ converge strongly, so (5) makes $u_k$ converge strongly to a unit vector in both $N(I-K)$ and its orthogonal complement, a contradiction.

For $v_k=(I-K)u_k\to v$, choose $u_k\in N(I-K)^\perp$. Estimate (4) gives
\[
\|v_k-v_l\|\geq\gamma\|u_k-u_l\|,
\]
so $u_k\to u$ and $(I-K)u=v$. Thus (ii) holds. Combining closedness with
\[
\overline{R(A)}=N(A^*)^\perp
\]
gives (iii).

If $N(I-K)=\{0\}$ but $R(I-K)\ne H$, the closed subspaces $H_j=(I-K)^jH$ form a strictly decreasing chain. Unit vectors $u_j\in H_j\cap H_{j+1}^\perp$ make $\{Ku_j\}$ uniformly separated, contradicting compactness. Hence injectivity implies surjectivity. Conversely, surjectivity and (iii) give $N(I-K^*)=\{0\}$; compactness of $K^*$ and the first implication give $R(I-K^*)=H$, and (iii) applied to $K^*$ gives $N(I-K)=\{0\}$. This proves (iv).

Finally, suppose $\dim N(I-K)<\dim R(I-K)^\perp$. Choose a finite-rank injection $A:N(I-K)\to R(I-K)^\perp$ that is not onto, and extend it by zero on $N(I-K)^\perp$. If $(I-K-A)u=0$, then $(I-K)u=Au$ lies in both $R(I-K)$ and $R(I-K)^\perp$; hence both sides vanish, and injectivity of $A$ on $N(I-K)$ gives $u=0$. Thus $K+A$ is compact with $N(I-K-A)=\{0\}$, so (iv) makes $I-K-A$ onto. Choosing $v\in R(I-K)^\perp\setminus R(A)$ contradicts this surjectivity. Therefore
\[
\dim N(I-K)\geq\dim R(I-K)^\perp.
\]
For $K$ this and (iii) give
\[
\dim N(I-K)\geq\dim N(I-K^*).
\]
Applying the same inequality to $K^*$ gives the reverse inequality, and therefore (v).
\end{proof}

\begin{remark}[The Fredholm alternative dichotomy]
\label{rem:fredholm-alternative-dichotomy}
\uses{thm:fredholm-alternative}
\dcref{appD:D.5-rem-2}
\group{Fredholm Theory}

Exactly one of the following alternatives holds:
\[
(\alpha)\qquad
\text{for every }f\in H,\ u-Ku=f
\text{ has a unique solution},
\]
or
\[
(\beta)\qquad
u-Ku=0\text{ has a nonzero solution}.
\]
In case $(\beta)$, the homogeneous solution space is finite dimensional, and
\[
(\gamma)\qquad u-Ku=f
\]
is solvable if and only if $f\in N(I-K^*)^\perp$.
\end{remark}

\begin{theorem}[Spectrum of a compact operator]
\label{thm:spectrum-of-compact-operator}
\uses{thm:fredholm-alternative, def:resolvent-set-spectrum-appendix-d, def:eigenvalue-eigenvector-appendix-d, def:compact-linear-operator}
\dcref{appD:thm-6}
\group{Fredholm Theory}

Assume $\dim H=\infty$ and $K:H\to H$ is compact. Then:

(i) $0\in\sigma(K)$;

(ii) $\sigma(K)-\{0\}=\sigma_p(K)-\{0\}$;

(iii) $\sigma(K)-\{0\}$ is finite, or it is a sequence tending to $0$.
\end{theorem}
\begin{proof}
\sketch If $0\notin\sigma(K)$, then $K^{-1}$ is bounded and $I=K\circ K^{-1}$ is compact, impossible on infinite-dimensional $H$. This proves (i). If $\eta\in\sigma(K)\setminus\{0\}$ and $N(K-\eta I)=\{0\}$, apply the Fredholm alternative to the compact operator $K/\eta$ to obtain $R(K-\eta I)=H$, contradicting $\eta\in\sigma(K)$. Hence (ii).

For distinct $\eta_k\in\sigma(K)\setminus\{0\}$ with $\eta_k\to\eta$, choose eigenvectors $Kw_k=\eta_kw_k$ and let $H_k=\operatorname{span}\{w_1,\ldots,w_k\}$. The eigenvectors are linearly independent and $(K-\eta_kI)H_k\subset H_{k-1}$. Choose unit vectors $u_k\in H_k\cap H_{k-1}^\perp$. For $k>l$,
\[
\frac{Ku_k}{\eta_k}-\frac{Ku_l}{\eta_l}
=\frac{Ku_k-\eta_ku_k}{\eta_k}
-\frac{Ku_l-\eta_lu_l}{\eta_l}
+u_k-u_l,
\]
and orthogonality gives
\[
\left\|\frac{Ku_k}{\eta_k}-\frac{Ku_l}{\eta_l}\right\|\geq1.
\]
If $\eta\ne0$, compactness of $K$ would make a subsequence of these ratios converge, a contradiction. Thus every accumulation point of the nonzero spectrum is $0$, proving (iii).
\end{proof}

\subsection{Symmetric Operators}
\label{sec:symmetric-operators}

Let $S:H\to H$ be bounded and symmetric, and set
\[
m:=\inf_{\substack{u\in H\\\|u\|=1}}(Su,u),
\qquad
M:=\sup_{\substack{u\in H\\\|u\|=1}}(Su,u).
\]

\begin{lemma}[Bounds on the spectrum of a symmetric operator]
\label{lem:bounds-on-spectrum-symmetric-operator}
\uses{def:adjoint-symmetric-operator, def:resolvent-set-spectrum-appendix-d, thm:lax-milgram-theorem, lem:cauchy-schwarz-inequality}
\dcref{appD:D.6-lem-1}
\group{Fredholm Theory}

(i) $\sigma(S)\subset[m,M]$;

(ii) $m,M\in\sigma(S)$.
\end{lemma}
\begin{proof}
\sketch If $\eta>M$, then
\[
((\eta I-S)u,u)\geq(\eta-M)\|u\|^2.
\]
The Lax--Milgram theorem (\cref{thm:lax-milgram-theorem}) makes $\eta I-S$ bijective, so $\eta\in\rho(S)$. Applying the same argument to $S-\eta I$ when $\eta<m$ proves (i).

For the upper endpoint, the symmetric form $[u,v]=((MI-S)u,v)$ is nonnegative. Cauchy--Schwarz (\cref{lem:cauchy-schwarz-inequality}) for this form gives
\[
(6)\qquad
\|(MI-S)u\|
\leq C((MI-S)u,u)^{1/2}.
\]
Choose unit vectors $u_k$ with $(Su_k,u_k)\to M$. Then (6) gives $(MI-S)u_k\to0$. If $M\in\rho(S)$, boundedness of $(MI-S)^{-1}$ would force $u_k\to0$, contradicting $\|u_k\|=1$. Thus $M\in\sigma(S)$. Applying the same argument to the nonnegative form $((S-mI)u,v)$ proves $m\in\sigma(S)$.
\end{proof}
